\subsubsection*{Der ISBN Code (ISBN10)}


\begin{definition}[Wort]
Ein \textbf{Wort} $\boldsymbol{w}$ ist ein Zeilenvektor aus $\Zm{11}^{10}$ bestehend aus 10 Elementen von $\Zm{11}$.
\[ \boldsymbol{w} = w_1\, w_2 \ldots w_{10} \]
\[ w_i \in \Zm{11} \quad \text{Buchstabe} \]
\end{definition}

$\Zm{11}^{10}$ ist ein Vektorraum über dem Körper $\Zm{11}$.\\
denn: Man kann Wörter addieren und mit Elementen von $\Zm{11}$ multiplizieren.

\begin{beispiel}
\[ \boldsymbol{\Zm{2}^3} = \{ \begin{array}[t]{llll} 000 & 001 & 010 & 011 \\ 100 & 101 & 110 & 111 \,\} \end{array} \]
\end{beispiel}

\begin{satz}
Die Anzahl der Wörter ist $|\Zm{11}^{10}| = 11^{10}$.
\end{satz}

Sei $C$ die Menge der \textbf{ISBN10}-Nummern (Menge der Codewörter).
\[ C \subseteq \Zm{11}^{10} \]

\begin{beispiel}
\textbf{0-387-98999-4}, \quad \textbf{3-528-06580-$\boldsymbol{X}$}
\end{beispiel}

\[ \boldsymbol{c} = c_1\text{-}c_2\, c_3\, c_4\text{-}c_5\, c_6\, c_7\, c_8\, c_9\text{-}c_{10} \in C \]
\begin{tabular}{ll}
$c_1$ & Sprache, \\
$c_2\, c_3\, c_4$ & Verlag, \\
$c_5 \ldots c_9$ & vergibt der Verlag, \\
$c_{10}$ & Prüfziffer.
\end{tabular}\\
$0 \leq c_1, \ldots, c_9 < 10$

Die Prüfziffer $c_{10}$ wird so bestimmt, dass die gewichtete Prüfsumme
\[ \sum_{i=1}^{10} i\, c_i \equiv 0 \pmod{11} \qquad (*) \]
verschwindet. Dann ist
\[ \sum_{i=1}^{10} i\, c_i = \sum_{i=1}^{9} i\, c_i + 10\, c_{10} \equiv \sum_{i=1}^{9} i\, c_i - c_{10} \equiv 0 \pmod{11} \]
\[ \Rightarrow \quad c_{10} \equiv \sum_{i=1}^{9} i\, c_i \pmod{11} \]
Man schreibt $X$ an Stelle von 10.

Sei $\boldsymbol{c} \in C$, $\boldsymbol{r} \in \Zm{11}^{10}$, $\boldsymbol{h} = 1\ 2\ 3\ 4\ 5\ 6\ 7\ 8\ 9\ 10$. Dann erhält $(*)$ die Form
\[ \boldsymbol{h}\, \boldsymbol{c}^T \equiv 0 \pmod{11}. \]
Wenn sich $\boldsymbol{r}$ und $\boldsymbol{c}$ an genau einer (Fehler-)Stelle $x$ unterscheiden, ist der Fehlervektor
\[ \boldsymbol{e} = \boldsymbol{r} - \boldsymbol{c} = 0\ 0 \ldots y \ldots 0\ 0, \qquad y = r_x - c_x \neq 0 \quad \text{Fehlergröße} \]

\begin{satz}
% KORRIGIERT: stand „erkennt genau einen Fehler“
ISBN10 erkennt jeden einzelnen Fehler (Fehler an genau einer Stelle).
\end{satz}

\begin{proof}
% KORRIGIERT: Malpunkt, „≢ … (mod 11)“ statt „≠“ und Schluss (Fehler wird erkannt) ergänzt
\[ \sum_{i=1}^{10} i\, r_i = \boldsymbol{h}\, \boldsymbol{r}^T = \boldsymbol{h}\, \boldsymbol{c}^T + \boldsymbol{h}\, \boldsymbol{e}^T \equiv \boldsymbol{h}\, \boldsymbol{e}^T = x \cdot y \pmod{11}, \quad x, y \not\equiv 0 \pmod{11} \]
Da $\Zm{11}$ ein Körper ist, gibt es keine Nullteiler, also
\[ x \cdot y \not\equiv 0 \pmod{11} \quad \Leftrightarrow \quad x \not\equiv 0 \text{ und } y \not\equiv 0 \pmod{11} \]
Also ist $\boldsymbol{h}\, \boldsymbol{r}^T \not\equiv 0 \pmod{11}$, d.h.\ $\boldsymbol{r} \notin C$: der Fehler wird erkannt.
\end{proof}

\begin{beispiel}
Sind obige Nummern ISBN10 Nummern?
\end{beispiel}
